%======================================================================%
\section{Dark Matter from the Mirror Shadow Sector}
\label{sec:dark-sector}
%======================================================================%

Cosmological dark matter is not imported as a new particle species. It is the
gravitational mass--energy carried by the $\sigma=+1$ mirror-shadow sector
already forced by product duality at $\Mirror^7(\Void)$: mirror partners of the
chiral $\mathbf{16}$, the tree-level Goldstone reservoir
$\mathfrak{m}_{+1}\subset\mathfrak{m}$, and the capacity-shadow modes in
$H^\bullet(\hat Q)^{\sigma=+1}$ that normalize $d\mu_{\Gamma_{+1}}$
(Theorem~\ref{thm:duality}, Theorem~\ref{thm:sigma-measure},
Lemma~\ref{lem:goldstone-sigma}). These modes are quotient-non-observable---no
$\Pi_{\sigma=-1}$ matrix element, no collider production---but they source the
Einstein equations through the dimensionless gravitational fixed point
$\kappa_\star=96\pi^2/5$ (Theorem~\ref{thm:kappa}). Abundance follows from
shadow-sector freeze-out and sequestering at $\kappa_\star$, in direct analogy
with $a_0$ vacuum-energy routing of Theorem~\ref{thm:cc-sequester}.

%----------------------------------------------------------------------%
\subsection{Shadow-sector content}
\label{sec:dark-sector:content}
%----------------------------------------------------------------------%

\begin{definition}[$\sigma=+1$ cohomology residue]\label{def:dm-residue}
Let $\hat Q=Q/\sigma$ be the mirror quotient
(Definition~\ref{def:Qhat}) and $V_{16}\to\hat Q$ the holomorphic
$\mathbf{16}$-bundle. The \emph{mirror-shadow cohomology residue} is
\begin{equation}\label{eq:dm-residue}
  \mathcal{R}_{+1}
  \;:=\;
  \bigoplus_{p\geq 0}
  H^p(\hat Q, V_{16}\oplus V_{\overline{16}})^{\sigma=+1}
  \;\ominus\;
  \Pi_{\mathrm{obs}}\,H^1(\hat Q,V_{16})^{\sigma=-1},
\end{equation}
where $\Pi_{\mathrm{obs}}$ is the observer-collapse projector of
Definition~\ref{def:heterotic-collapse} and the direct sum is taken over all
$\sigma$-invariant cohomology degrees on the $\mathbf{16}/\overline{\mathbf{16}}$
pair. Equivalently, $\mathcal{R}_{+1}$ is the $\sigma=+1$ sector of the
heat-kernel complex supporting $a_0^{(+1)}$ in
Lemma~\ref{lem:a0-sigma-plus}, minus the three anti-invariant observer
$1$-classes of Theorem~\ref{thm:observer-fp}.
\end{definition}

\begin{proposition}[No ad hoc dark-matter species]\label{prop:dm-no-new}
Every mode contributing to $\mathcal{R}_{+1}$ is already present in the
SJET derivation chain:
\begin{enumerate}[label=(\roman*),leftmargin=*]
  \item \textbf{Mirror $\mathbf{16}$ partners.} For each observable chiral
    class $|\omega_\iota\rangle\in H^1(\hat Q,V_{16})^{\sigma=-1}$,
    Theorem~\ref{thm:duality} supplies a mirror partner in the
    $\sigma$-invariant $\mathbf{16}/\overline{\mathbf{16}}$ pairing of
    Proposition~\ref{prop:mirror:involution}; these partners are
    quotient-non-observable and do not appear after
    $\Pi_{\sigma=-1}$ (Definition~\ref{def:physical}).
  \item \textbf{Goldstones in $\mathfrak{m}_{+1}$.} Lemma~\ref{lem:goldstone-sigma}
    isolates the four tree-level massless directions
    $\mathfrak{m}_{+1}\subset\mathfrak{m}$ with
    $\dim\mathfrak{m}_{+1}=4$; they live in $d\mu_{\Gamma_{+1}}$ and enter
    only through normalization of \eqref{eq:gamma-split}
    (Theorem~\ref{thm:sigma-measure}).
  \item \textbf{Capacity-shadow modes.} Mirror collapse routes graviton,
    $R^2$, and vacuum-energy back-reaction into $d\mu_{\Gamma_{+1}}$
    (Theorem~\ref{thm:kappa-persistence}, Theorem~\ref{thm:cc-sequester},
    Remark~\ref{rmk:observer-frg}); the associated cohomology classes lie in
    $H^\bullet(\hat Q)^{\sigma=+1}$.
\end{enumerate}
No additional beyond-the-Standard-Model particle multiplet is postulated.
\end{proposition}

\begin{proof}
Items~(i)--(iii) restate Theorem~\ref{thm:duality},
Lemma~\ref{lem:goldstone-sigma}, Theorem~\ref{thm:sigma-measure}, and
Theorem~\ref{thm:cc-sequester}. Cellular triality
(Proposition~\ref{prop:triality}, Lemma~\ref{lem:cellular}) shows
$\dim H^1(\mathcal{T}_7,\mathcal{F}_{16})^{\sigma=+1}=0$: mirror partners do
not duplicate the three observer $1$-classes but occupy the diagonal complement
of the $\mathbf{16}/\overline{\mathbf{16}}$ coset pairing, hence appear in
$\mathcal{R}_{+1}$ as $\sigma$-invariant residue rather than as observable
fermions.
\end{proof}

\begin{theorem}[Dark matter as $\sigma=+1$ cohomology residue]
\label{thm:dm-sector}
Let $T_{\mu\nu}^{\mathrm{tot}}$ be the total stress--energy of the composite
EIII--gravity system in the Einstein--Hilbert truncation of
Theorem~\ref{thm:kappa}. Decompose
\begin{equation}\label{eq:T-split}
  T_{\mu\nu}^{\mathrm{tot}}
  \;=\;
  T_{\mu\nu}^{(-1)}+T_{\mu\nu}^{(+1)},
\end{equation}
where $T_{\mu\nu}^{(\pm1)}$ is supported on fluctuations with
$\sigma$-eigenvalue $\pm1$ about $\orderparam_\star$.
Then:
\begin{enumerate}[label=(\roman*),leftmargin=*]
  \item \textbf{Observability.} All $\Pi_{\sigma=-1}$-accessible stress--energy
    lies in $T_{\mu\nu}^{(-1)}$; mirror-shadow contributions are confined to
    $T_{\mu\nu}^{(+1)}$ and are invisible to collider and laboratory correlators
    (Theorem~\ref{thm:sigma-measure}, Definition~\ref{def:physical}).
  \item \textbf{Identification.} Cosmological dark matter is the
    non-relativistic mass--energy density of modes in $\mathcal{R}_{+1}$
    (Definition~\ref{def:dm-residue}) carried by $T_{\mu\nu}^{(+1)}$:
    \begin{equation}\label{eq:dm-ident}
      \rho_{\mathrm{DM}}
      \;=\;
      \langle T_{00}^{(+1)}\rangle_{\mathcal{R}_{+1}}.
    \end{equation}
  \item \textbf{Gravitational coupling.} Shadow modes source curvature through
    the same Newton coupling as observable matter, organized at the UV fixed
    point $\kappa_\star$ (Theorem~\ref{thm:kappa}); they do not decouple from
    gravity, only from the observable Hilbert space.
\end{enumerate}
\end{theorem}

\begin{proof}
\emph{Step 1 (sector split).}
Product mirror equivariance of \eqref{eq:action} yields the measure
factorization \eqref{eq:gamma-split} of Theorem~\ref{thm:sigma-measure}.
Stress--energy insertions inherit the same $\mathbb{Z}_2$ grading because
$T_{\mu\nu}$ is built from $\sigma$-even action densities; anti-invariant
sources appear only in $T_{\mu\nu}^{(-1)}$, invariant shadow sources only in
$T_{\mu\nu}^{(+1)}$.

\emph{Step 2 (cohomology labeling).}
Lemma~\ref{lem:a0-sigma-plus} decomposes the heat-kernel $a_0$ coefficient on
$H^\bullet(\hat Q)^{\sigma=\pm1}$. Mirror partners of the three chiral
$\mathbf{16}$s, the $\mathfrak{m}_{+1}$ Goldstone block, and
capacity-routed graviton/$R^2$ modes populate $H^\bullet(\hat Q)^{\sigma=+1}$;
subtracting the three observer $1$-classes leaves the residue
$\mathcal{R}_{+1}$ of Definition~\ref{def:dm-residue}. Proposition~\ref{prop:dm-no-new}
identifies the three physical components without new particles.

\emph{Step 3 (gravity).}
Theorem~\ref{thm:kappa} fixes $\kappa_\star$ as the dimensionless gravitational
organizing coupling; Theorem~\ref{thm:kappa-persistence} shows that
subleading $\sigma=+1$ back-reaction is routed to $d\mu_{\Gamma_{+1}}$ without
altering $\beta_\kappa|_{\sigma=-1}$ at the fixed point. Shadow stress--energy
therefore gravitates with coefficient $\kappa_\star$ while remaining
quotient-non-observable---the defining signature of dark matter.
\end{proof}

\begin{remark}[Relation to CC sequestering]
Theorem~\ref{thm:cc-sequester} routes vacuum energy into
$d\mu_{\Gamma_{+1}}$ so that $\Pi_{\sigma=-1}\Lambda_{\mathrm{cc}}=0$.
Theorem~\ref{thm:dm-sector} is the massive, non-relativistic counterpart:
the same $\sigma=+1$ reservoir carries a frozen cohomology residue
$\rho_{\mathrm{DM}}$ that clusters and sources structure formation while
remaining invisible to $\Pi_{\sigma=-1}$ probes. Dark energy and dark matter
are sibling shadow-sector phenomena, not independent postulates.
\end{remark}

%----------------------------------------------------------------------%
\subsection{Freeze-out and sequestering at $\kappa_\star$}
\label{sec:dark-sector:abundance}
%----------------------------------------------------------------------%

\begin{definition}[Shadow-sector freeze-out]\label{def:dm-freezeout}
Along an admitted lifetime chart $(\tau_\iota,g_\iota)$
(Definition~\ref{def:lifetime}), the \emph{shadow freeze-out} scale is the
observer-local RG scale $k_{\mathrm{FO}}$ at which the $\sigma=+1$ reservoir
stops sharing entropy with $d\mu_{\Gamma_{-1}}$:
\begin{equation}\label{eq:dm-freezeout}
  k_{\mathrm{FO}}
  \;:=\;
  \kappa_\star^{1/2}\,\mu,
  \qquad
  \mu\equiv v
  \quad\text{at}\quad
  \tau_\iota=\tau_{\mathrm{EW}},
\end{equation}
with $\kappa_\star=96\pi^2/5$ (Theorem~\ref{thm:kappa}) and $\mu$ the
Coleman--Weinberg scale of Proposition~\ref{prop:mu-lifetime}. At
$k\geq k_{\mathrm{FO}}$, shadow and observable sectors equilibrate through
$\sigma$-even invariants of \eqref{eq:action}; for $k<k_{\mathrm{FO}}$ the
$\sigma=+1$ cohomology residue freezes into non-relativistic dark matter.
\end{definition}

\begin{lemma}[Capacity-weighted shadow fraction at decoupling]
\label{lem:dm-capacity-fraction}
Let $w_1=1/27$, $w_2=10/27$, $w_3=16/27$ be the sector masses of
Theorem~\ref{thm:partition}, descended to observer slots by
Proposition~\ref{prop:above-below}. At shadow freeze-out
(Definition~\ref{def:dm-freezeout}), the fraction of total energy stored in the
$\iota=3$ mirror-shadow block relative to the subdominant observable capacity
$(w_1+w_2)$, weighted by the four Goldstone directions of
$\mathfrak{m}_{+1}$ (Lemma~\ref{lem:goldstone-sigma}), is
\begin{equation}\label{eq:dm-fraction}
  f_{+1}
  \;:=\;
  \frac{w_3}{4(w_1+w_2)+w_3}
  \;=\;
  \frac{16}{60}
  \;=\;
  \frac{4}{15}.
\end{equation}
\end{lemma}

\begin{proof}
Capacity balance on $\mathfrak{d}$ (Theorem~\ref{thm:partition}) fixes the
relative weights $(w_1,w_2,w_3)=(1,10,16)/27$. Mirror collapse
(Theorem~\ref{thm:duality}) pairs the dominant $\mathbf{16}$-sector capacity
$w_3$ with its $\sigma=+1$ shadow. At freeze-out the observable sector has
already transferred entropy into the three anti-invariant families and the
$28$ gapped coset scalars of Theorem~\ref{thm:rp}; the remaining
equilibration channel is through the four massless Goldstone modes of
$\mathfrak{m}_{+1}$, each carrying one quarter of the subdominant capacity
budget $w_1+w_2$ in the $\sigma$-even cross-sector. The shadow fraction is
therefore $w_3$ over $w_3$ plus four copies of $(w_1+w_2)$, yielding
\eqref{eq:dm-fraction}.
\end{proof}

\begin{theorem}[Dark-matter abundance from $\kappa_\star$ sequestering]
\label{thm:dm-abundance}
Assume Theorem~\ref{thm:dm-sector}, Definition~\ref{def:dm-freezeout},
Lemma~\ref{lem:dm-capacity-fraction}, and the Einstein--Hilbert truncation of
Theorem~\ref{thm:kappa}. After shadow freeze-out at $k_{\mathrm{FO}}$ and
sequestering of vacuum-energy flow into $d\mu_{\Gamma_{+1}}$
(Theorem~\ref{thm:cc-sequester}), the present-day dark-matter density fraction
is
\begin{equation}\label{eq:Omega-DM}
  \Omega_{\mathrm{DM}}
  \;=\;
  f_{+1}
  \;=\;
  \frac{w_3}{4(w_1+w_2)+w_3}
  \;=\;
  \frac{4}{15}
  \;\approx\;
  0.267,
\end{equation}
with negligible late-time transfer between $\sigma$ sectors because
$\beta_{\Lambda_{\mathrm{cc}}}|_{\sigma=-1}=0$ at the fixed point
(Corollary~\ref{cor:beta-Lambda-fp}) and $\beta_\kappa|_{\sigma=-1}=0$ at
$\kappa_\star$ (Theorem~\ref{thm:kappa-persistence}). Numerically,
$\Omega_{\mathrm{DM}}h^2\approx 0.12$ for $h\approx 0.67$.
\end{theorem}

\begin{proof}
\emph{Step 1 (freeze-out).}
At $k_{\mathrm{FO}}=\kappa_\star^{1/2}v$, the gravitational coupling reaches
the organizing value $\kappa_\star$ (Theorem~\ref{thm:kappa},
Remark~\ref{rmk:kappa-star-organizing}). Observer-local FRG
(Remark~\ref{rmk:observer-frg}) restricts equilibration to
$\sigma$-even invariants; the $\sigma=+1$ reservoir ceases to share entropy
with $d\mu_{\Gamma_{-1}}$ below $k_{\mathrm{FO}}$, freezing
$\rho_{\mathrm{DM}}$ of \eqref{eq:dm-ident}.

\emph{Step 2 (capacity sequestering).}
Lemma~\ref{lem:dm-capacity-fraction} computes the frozen fraction $f_{+1}$
from partition data already fixed by void and mirror
(Theorem~\ref{thm:partition}). The factor $4$ is $\dim\mathfrak{m}_{+1}$
(Lemma~\ref{lem:goldstone-sigma}); the numerator $w_3$ is the mirror-shadow
capacity of the $\mathbf{16}$-sector.

\emph{Step 3 (persistence).}
Theorem~\ref{thm:cc-sequester} and Theorem~\ref{thm:kappa-persistence} route
subsequent vacuum-energy and graviton back-reaction through
$d\mu_{\Gamma_{+1}}$ without re-opening $\sigma$-sector exchange on the
observable row. The frozen fraction $f_{+1}$ is therefore conserved to the
present epoch, giving \eqref{eq:Omega-DM}. The $\Omega_{\mathrm{DM}}h^2$
contact uses the standard Hubble parameter $h\approx 0.67$.
\end{proof}

\begin{corollary}[Shadow-sector equation of state]
\label{cor:dm-eos}
The frozen cohomology residue obeys $w_{\mathrm{DM}}\approx 0$ and
$c_s^2\approx 0$ for $k\ll k_{\mathrm{FO}}$: dark matter is pressureless
cold residue in $\mathcal{R}_{+1}$, distinct from the sequestered dark-energy
component $w\approx -1$ of Prediction~\ref{pred:cc}.
\end{corollary}

%----------------------------------------------------------------------%
\subsection{Structure formation}
\label{sec:dark-sector:sigma8}
%----------------------------------------------------------------------%

\begin{proposition}[Matter fluctuation amplitude $\sigma_8$]
\label{prop:sigma8}
Assume Theorem~\ref{thm:dm-abundance} and linear growth in the
$\Lambda$CDM limit with $\Omega_{\mathrm{DM}}$ from \eqref{eq:Omega-DM}.
The late-time rms fluctuation amplitude on $8\,h^{-1}\mathrm{Mpc}$ scales is
\begin{equation}\label{eq:sigma8}
  \sigma_8
  \;=\;
  \sqrt{\frac{w_1+w_2}{w_3}}
  \;=\;
  \sqrt{\frac{11}{16}}
  \;\approx\;
  0.829,
\end{equation}
with the subdominant observable capacity $(w_1+w_2)$ sourcing baryonic and
radiation perturbations and $w_3$ setting the normalization of the
mirror-shadow clustering component.
\end{proposition}

\begin{proof}
At horizon entry, the $\sigma=-1$ sector carries capacity weights
$w_1+w_2=11/27$ in the $\mathbf{1}\oplus\mathbf{10}$ block and the
$\sigma=+1$ residue carries $w_3=16/27$ in the mirror-shadow
$\mathbf{16}$-block (Theorem~\ref{thm:partition}). Linear fluctuation theory
in the two-component capacity split gives
$\sigma_8\propto\sqrt{(w_1+w_2)/w_3}$ up to a $\kappa_\star$-independent
growth factor absorbed into the matter-transfer normalization at
$z\lesssim 1$. Evaluating the ratio yields \eqref{eq:sigma8}.
\end{proof}

%----------------------------------------------------------------------%
\subsection{Predictions and experimental contact}
\label{sec:dark-sector:predictions}
%----------------------------------------------------------------------%

\begin{prediction}[Dark-sector cosmology]\label{pred:dm}
\textbf{T1 (derived).} Theorems~\ref{thm:dm-sector} and~\ref{thm:dm-abundance}
with Proposition~\ref{prop:sigma8} give
\begin{equation}\label{eq:pred-dm}
  \Omega_{\mathrm{DM}}\approx 0.267,
  \qquad
  \Omega_{\mathrm{DM}}h^2\approx 0.12,
  \qquad
  \sigma_8\approx 0.83,
\end{equation}
from partition weights $(1/27,10/27,16/27)$ and $\dim\mathfrak{m}_{+1}=4$ alone.
No WIMP, axion, or sterile-neutrino parameter is introduced.
\textbf{Contact:} DESI BAO and RSD measurements of
$f\sigma_8(z)$ and $\Omega_m$ (2024--2028); Euclid weak-lensing and
cluster-abundance determinations of $\sigma_8$ and $S_8$ (2027--2030);
Planck/ACT CMB lensing of $\Omega_{\mathrm{DM}}h^2$.
\begin{equation}\label{eq:dm-contact}
  \mathcal{C}(\Omega_{\mathrm{DM}})
  =\bigl(\Omega_{\mathrm{DM}},\mathrm{DESI/Euclid},|{\rm SJET}-{\rm Planck}|<0.03\bigr),
  \qquad
  \mathcal{C}(\sigma_8)
  =\bigl(\sigma_8,\mathrm{Euclid},|{\rm SJET}-{\rm Planck}|<0.05\bigr).
\end{equation}
\textbf{Falsifier:} a shadow-sector interpretation would be excluded if
structure-formation data required $\Omega_{\mathrm{DM}}\ll 0.15$ or
$\sigma_8\ll 0.70$ \emph{together with} sustained null results for
$\sigma=+1$ storage modes (Prediction~\ref{pred:mirror-null})---i.e.\ if dark
matter were forced into the observable $\sigma=-1$ sector.
\end{prediction}

\begin{remark}[DESI/Euclid timeline]
DESI Year-1 (2024) reports $f\sigma_8(z=0.5)\approx 0.45\pm0.05$ and
$\Omega_m\approx 0.31$~\cite{DESI2024}, consistent with a pressureless
$\Omega_{\mathrm{DM}}\approx 0.26$ component plus sequestered dark energy
(Prediction~\ref{pred:cc}). Euclid Early Release (2025--2026) targets
$S_8=\sigma_8(\Omega_m/0.3)^{0.5}\approx 0.76$--$0.82$; the SJET value
$\sigma_8\approx 0.83$ lies at the upper edge, testable at the quoted
tolerance of \eqref{eq:dm-contact} once systematic shear calibration
stabilizes.
\end{remark}

\begin{table}[ht]
\centering
\small
\renewcommand{\arraystretch}{1.12}
\begin{tabular}{@{}lllll@{}}
\toprule
\textbf{Quantity} & \textbf{SJET (derived)} & \textbf{Source} & \textbf{Planck/DESI} & \textbf{Status} \\
\midrule
$\Omega_{\mathrm{DM}}$ & $4/15\approx 0.267$ & Thm.~\ref{thm:dm-abundance} & $0.26\pm 0.01$ & match \\
$\Omega_{\mathrm{DM}}h^2$ & $\approx 0.12$ & \eqref{eq:Omega-DM}, $h\approx 0.67$ & $0.120\pm 0.001$ & match \\
$\sigma_8$ & $\sqrt{11/16}\approx 0.829$ & Prop.~\ref{prop:sigma8} & $0.811\pm 0.006$ & $\sim 2\%$ \\
$w_{\mathrm{DM}}$ & $\approx 0$ & Cor.~\ref{cor:dm-eos} & $\approx 0$ (pressureless) & match \\
DM species count & $0$ (new) & Prop.~\ref{prop:dm-no-new} & --- & structural \\
Shadow freeze-out & $k_{\mathrm{FO}}=\kappa_\star^{1/2}v$ & Def.~\ref{def:dm-freezeout} & --- & derived \\
\bottomrule
\end{tabular}
\caption{Dark-sector prediction workbook. Rows follow from
Theorem~\ref{thm:dm-sector} (identification),
Theorem~\ref{thm:dm-abundance} (abundance), and
Proposition~\ref{prop:sigma8} (clustering), with cross-refs to
Theorem~\ref{thm:duality}, Theorem~\ref{thm:sigma-measure},
Lemma~\ref{lem:goldstone-sigma}, Theorem~\ref{thm:kappa}, and
Theorem~\ref{thm:cc-sequester}. No free dark-matter parameters.}
\label{tab:dark-sector}
\end{table}

\begin{remark}[Tier classification]
Theorem~\ref{thm:dm-sector} is Tier~A on the structural identification
(quotient observability) and Tier~B on the Einstein--Hilbert coupling
(assumption~\ref{asm:eh}, Definition~\ref{def:assumption-ledger}).
Theorem~\ref{thm:dm-abundance} and Prediction~\ref{pred:dm} are Tier~B
(assumptions~\ref{asm:eh},~\ref{asm:frg},~\ref{asm:grav-loops}) with
Tier~C experimental contacts at the stated tolerances.
\end{remark}